%% Bearbeite nur diese Datei und sample-base.bib. 
%% dann passiert nix 
%% This is file `sample-acmtog.tex',
%% generated with the docstrip utility.
%%
%% The original source files were:
%%
%% samples.dtx  (with options: `all,journal,bibtex,acmtog')
%% 
%% IMPORTANT NOTICE:
%% 
%% For the copyright see the source file.
%% 
%% Any modified versions of this file must be renamed
%% with new filenames distinct from sample-acmtog.tex.
%% 
%% For distribution of the original source see the terms
%% for copying and modification in the file samples.dtx.
%% 
%% This generated file may be distributed as long as the
%% original source files, as listed above, are part of the
%% same distribution. (The sources need not necessarily be
%% in the same archive or directory.)
%%
%%
%% Commands for TeXCount
%TC:macro \cite [option:text,text]
%TC:macro \citep [option:text,text]
%TC:macro \citet [option:text,text]
%TC:envir table 0 1
%TC:envir table* 0 1
%TC:envir tabular [ignore] word
%TC:envir displaymath 0 word
%TC:envir math 0 word
%TC:envir comment 0 0
%%
%%
%% The first command in your LaTeX source must be the \documentclass
%% command.
%%
%% For submission and review of your manuscript please change the
%% command to \documentclass[manuscript, screen, review]{acmart}.
%%
%% When submitting camera ready or to TAPS, please change the command
%% to \documentclass[sigconf]{acmart} or whichever template is required
%% for your publication.
%%
%%
\documentclass[acmtog, nonacm=true, language=german]{acmart}
\usepackage[normalem]{ulem} %%% Ermöglicht durchstreichen!
\usepackage{pdflscape} %ermöglicht querformat
%%
\usepackage{tabularx}
\usepackage{adjustbox} %ermöglicht tabellen format
\usepackage{lipsum}
%% \BibTeX command to typeset BibTeX logo in the docs
\AtBeginDocument{%
  \providecommand\BibTeX{{%
    Bib\TeX}}}

%% Rights management information.  This information is sent to you
%% when you complete the rights form.  These commands have SAMPLE
%% values in them; it is your responsibility as an author to replace
%% the commands and values with those provided to you when you
%% complete the rights form.
\setcopyright{none}
\copyrightyear{}
\acmYear{}
\acmDOI{}


%%
%% These commands are for a JOURNAL article.
\acmJournal{FACMP}
%\acmVolume{37}
%\acmNumber{4}
%\acmArticle{111}
%\acmMonth{8}

%%
%% Submission ID.
%% Use this when submitting an article to a sponsored event. You'll
%% receive a unique submission ID from the organizers
%% of the event, and this ID should be used as the parameter to this command.
%%\acmSubmissionID{123-A56-BU3}

%%
%% For managing citations, it is recommended to use bibliography
%% files in BibTeX format.
%%
%% You can then either use BibTeX with the ACM-Reference-Format style,
%% or BibLaTeX with the acmnumeric or acmauthoryear sytles, that include
%% support for advanced citation of software artefact from the
%% biblatex-software package, also separately available on CTAN.
%%
%% Look at the sample-*-biblatex.tex files for templates showcasing
%% the biblatex styles.
%%

%%
%% The majority of ACM publications use numbered citations and
%% references.  The command \citestyle{authoryear} switches to the
%% "author year" style.
%%
%% If you are preparing content for an event
%% sponsored by ACM SIGGRAPH, you must use the "author year" style of
%% citations and references.
\citestyle{acmauthoryear}
%\bibliographystyle{IEEEtran}


%%
%% end of the preamble, start of the body of the document source.
\begin{document}

%%
%% The "title" command has an optional parameter,
%% allowing the author to define a "short title" to be used in page headers.
\title{Bedrohungen und Gegenmaßnahmen im Entwicklungszyklus einer Künstlichen Intelligenz}

%%
%% The "author" command and its associated commands are used to define
%% the authors and their affiliations.
%% Of note is the shared affiliation of the first two authors, and the
%% "authornote" and "authornotemark" commands
%% used to denote shared contribution to the research.
\author{Max Mustermann}
%\authornote{Both authors contributed equally to this research.}
\email{mustermann@hm.edu}
%\orcid{1234-5678-9012}
\authornote{Max Mustermann (email@hm.edu): Wirtschaftsinformatikstudent seit 10/20YY. ... Weitere Infos ...}
%\author{G.K.M. Tobin}
%\authornotemark[1]
%\email{webmaster@marysville-ohio.com}
\affiliation{%
  \institution{Hochschule München}
  %\city{München}
  %\state{Bayern}
  \country{Deutschland}
}
\authorsaddresses{} %leer um es zu unterdrücken



%%
%% By default, the full list of authors will be used in the page
%% headers. Often, this list is too long, and will overlap
%% other information printed in the page headers. This command allows
%% the author to define a more concise list
%% of authors' names for this purpose.
\renewcommand{\shortauthors}{max}

%%
%% The abstract is a short summary of the work to be presented in the
%% article.
\begin{abstract}
\lipsum[1]
\end{abstract}
%%
%% The code below is generated by the tool at http://dl.acm.org/ccs.cfm.
%% Please copy and paste the code instead of the example below.
%%

%%
%% Keywords. The author(s) should pick words that accurately describe
%% the work being presented. Separate the keywords with commas.
\keywords{Wort1, Wort2, Wort3, Wort4}

%\received{20 February 2007}
%\received[revised]{12 March 2009}
%\received[accepted]{5 June 2009}

%%
%% This command processes the author and affiliation and title
%% information and builds the first part of the formatted document.
\maketitle

\section{Idee des Artikels}
\lipsum[1]
\section{Forschungsdesign}
\subsection{Forschungsfrage und Hypothese}
\lipsum[2]
\subsection{Datenerhebung}
\lipsum[3]
\subsection{Datenanalyse}
\lipsum[4]
\subsection{Ergebnisse}
\lipsum[5]
\subsection{Forschungsphylosophie}
\lipsum[6]
\section{State of the art}
\lipsum[7]
\section{Kritische Diskussion}
\subsection{Forschungsfrage und Hypothese}
\lipsum[8]
\subsection{Datenerhebung}
Hier ist ein Zitat von Adams \cite{adams_research_2007}. \lipsum[9]
\subsection{Datenanalyse}
\lipsum[10]
\subsection{Ergebnis}
\lipsum[11]
\subsection{Forschungsphylosophie}
    \begin{figure}
        \centering
        \includegraphics[width=\columnwidth]{saunders-onion.png}
        \caption{Research Onion nach Saunders}
        \label{fig:saundersOnion}
    \end{figure}

Die Research Onion nach Saunders \cite{saunders_research_2019} ist hier.
\lipsum[12]
\subsection{Praktische und wissenschaftliche Relevanz}
\lipsum[13]
\section{Fazit}
\subsection{Lessons Learned}
\lipsum[19]
\section{Literaturverzeichnis}

  \bibliographystyle{ACM-Reference-Format}
  \bibliography{sample-base}

\appendix 
\begin{landscape}

\section{Tabellarische Darstellung ''State of the art''}

\begin{table}[htbp]
\caption{Fachlich}
\label{tab:sota-fachlich}
% Setzen Sie eine passende Schriftgröße, um die Lesbarkeit zu gewährleisten.
%\footnotesize % oder \small, wenn notwendig
%umbruch mit \par
\begin{tabularx}{\linewidth}{|X|X|X|X|X|X|X|}
\hline
\textbf{Author/Jahr/Journal} & \textbf{Titel} & \textbf{Ziel der Forschungsarbeit} & \textbf{Datenerhebung} & \textbf{Datenanalyse} & \textbf{Forschungsergebnis} & \textbf{Bezug zum analysierten Artikel} \\
\hline
Hu, Y., Kuang, W., Qin, Z., Li, K., Zhang, J., Gao, Y., Li, W., \& Li, K./2021/ACM Computing Surveys &
  Artificial Intelligence Security: Threats and Countermeasures &
  Framework zu Gefahren und Risiken in unterschiedlichen Entwicklungsphasen einer KI herausarbeiten &
  Literatur &
  Lesen &
  Framework &
  100 \\
\hline
Hu, Y., Kuang, W., Qin, Z., Li, K., Zhang, J., Gao, Y., Li, W., \& Li, K./2021/ACM Computing Surveys &
  Artificial Intelligence Security: Threats and Countermeasures &
  Framework zu Gefahren und Risiken in unterschiedlichen Entwicklungsphasen einer KI herausarbeiten &
  Literatur &
  Lesen &
  Framework &
  100 \\
\hline
Hu, Y., Kuang, W., Qin, Z., Li, K., Zhang, J., Gao, Y., Li, W., \& Li, K./2021/ACM Computing Surveys &
  Artificial Intelligence Security: Threats and Countermeasures &
  Framework zu Gefahren und Risiken in unterschiedlichen Entwicklungsphasen einer KI herausarbeiten &
  Literatur &
  Lesen &
  Framework &
  100 \\
\hline
% Fügen Sie hier zusätzliche Zeilen ein
\end{tabularx}
\end{table}

\end{landscape}
\end{document}